\cotefile{45}
\scanpage{33}
\margin{32}
- \mspage{22} - \hfill n$^{\mathrm{o}}$ 246

la réunion de $]\Delta_1,\Delta_2[$, de $\Delta_1$ et de $\Delta_2$.

\uline{Lorsque} $k$ \uline{satisfait à l'hypothèse du chapeau du n$^{\mathrm{o}}$} 6, on peut, au moyen d'une forme quadratique auxiliaire, établir une correspondance biunivoque \uline{croissante} entre le corps $k$ et l'ensemble ordonné défini par la prop. 6. [N.B. Pour le cas où $k$ ne satisfait pas à l'hypothèse en question, l'adjudant de service fabriquera un contre-exemple et le mettra en exercice]. Pour cela, on introduit dans le plan $E$ n'importe quelle forme quadratique $\phi$ pour laquelle $-\delta$ est un carré, ce qui nous permet d'appliquer tous les résultats des n$^{\mathrm{os}}$ précédents. Pour toute demi-droite $\Delta$ autre que $\Delta_0$, on posera
\[
t(\Delta)=\operatorname{tg}\frac12(-\Delta_0,\Delta).
\]
Comme on a $\Delta\ne\Delta_0$, $(-\Delta_0,\Delta)$ n'est pas l'angle plat, donc $\frac12(-\Delta_0,\Delta)$ n'est pas un angle droit et par suite $t(\Delta)\ne\infty$. Il est immédiat que $\Delta\longrightarrow t(\Delta)$ est une bijection de l'ensemble des demi-droites $\ne\Delta_0$ sur $k$. On va montrer que cette bijection est croissante. Soient en effet $\Delta_1,\Delta_2$ deux demi-droites distinctes de $\Delta_0$ ; soit $t_i=t(\Delta_i)$ $(i=1,2)$
Soit $a$ un point, autre que $0$, sur $-\Delta_0$ ; pour la base $(a,Ja)$ de $E$, les demi-droites $\Delta_0,\Delta_1,\Delta_2$ contiendront respectivement (en vertu de la formule (12)) les points :
\[
\begin{aligned}
x_0&=(-1,0),\\
x_i&=(1-t_i^2,2t_i)\quad(i=1,2),
\end{aligned}
\]
Pour que la suite $\Delta_0,\Delta_1,\Delta_2$ soit directe, il faut et il suffit que deux au moins des quantités
\[
-2t_1,\quad(1-t_1^2)2t_2-(1-t_2^2)2t_1,\quad2t_2
\]
soient $>0$. Or la seconde s'écrit $2(t_2-t_1)(1+t_1t_2)$.
